\section{РЕЗУЛЬТАТЫ РАБОТЫ}

\subsection{Исходные данные}
\label{sec:final_data}

В качестве источника данных для графов дорожной сети выбраны три фрагмента карты Москвы разного размера.
После процедуры обработки получены следующие графы.

Малый граф состоит из 75 ребер (10 узких) и 68 вершин (1 базовая точка и 8 точек доставки) и 
представлен на рисунке~\ref{fig:small_graph} (точки доставки обозначены зеленым цветом, 
базовые точки --- красным, узкие ребра --- желтой волнистой линией).

\begin{figure}
\includegraphics[width=0.75\textwidth]{small_graph}
\caption{Малый граф}
\label{fig:small_graph}
\end{figure}

Средний граф состоит из 1054 ребер (102 узких) и 887 вершин (6 базовых точек и 73 точек доставки).

Большой граф состоит из 3469 ребер (250 узких) и 2871 вершин (17 базовых точек и 264 точки доставки).

Средний и большой графы представлены в приложении~\ref{app:graphs} на рисунках~\ref{fig:medium_graph} и~\ref{fig:large_graph}, соответственно.

\subsection{Используемые метрики}

Сравнение методов проводится по интегральным характеристикам потока заданий на фиксированном промежутке 
времени $[0, T]$, где $T$ --- длительность одного запуска симуляции. В ходе работы системы каждый агент 
$i \in A$~(\ref{e:agents}) многократно совершает циклы вида \enquote{базовая точка $\to$ точка доставки
$\to$ базовая точка}. Каждое перемещение агента по маршруту от одной из этих точек до следующей 
обозначим индексом $\ell$. Обозначим через $\mathcal{R}(T)$ множество маршрутов, завершившихся в пределах
$[0, T]$, а через $\mathcal{O}(T) \subseteq \mathcal{R}(T)$ --- подмножество маршрутов, завершающих 
доставку заказа. Для каждого маршрута $\ell$ известны его длительность $\tau_\ell$ и средняя скорость
движения $\bar v_\ell = L_\ell / \tau_\ell$, где $L_\ell$ --- пройденное агентом за маршрут расстояние.

В качестве основных показателей качества используются три величины:

\begin{itemize}
    \item число выполненных заказов за время наблюдения:

        \begin{equation}
            N(T) = |\mathcal{O}(T)|;
        \label{e:metric-orders}
        \end{equation}

    \item среднее время выполнения маршрута:

        \begin{equation}
            \bar\tau(T) = \frac{1}{|\mathcal{R}(T)|} \sum_{\ell \in \mathcal{R}(T)} \tau_\ell;
        \label{e:metric-order-time}
        \end{equation}

    \item средняя скорость движения агента по маршруту:

        \begin{equation}
            \bar v(T) = \frac{1}{|\mathcal{R}(T)|} \sum_{\ell \in \mathcal{R}(T)} \bar v_\ell.
        \label{e:metric-speed}
        \end{equation}
\end{itemize}

Метрика $N(T)$ обобщает на продолжительный режим работы число достигнутых целей из классической MAPF, 
$\bar\tau(T)$ --- метрику $\mathrm{SOC}$~(\ref{e:soc}), а $\bar v(T)$ характеризует эффективность 
исполнения индивидуальных маршрутов и снижается при возникновении заторов, ожиданий и отступлений.

% вспомогательные метрики
Дополнительно фиксируется ряд вспомогательных метрик, имеющих смысл не для всех
конфигураций и используемых для детального анализа работы отдельных механизмов:

\begin{itemize}
    \item $\bar\tau_{\mathrm{wait}}(T)$ --- среднее время непрерывного пребывания
        агента в состоянии ожидания;

    \item $\bar\tau_{\mathrm{rev}}(T)$ --- среднее время непрерывного пребывания агента в
        состоянии отступления;

    \item $N_c(T)$ --- суммарное число конфликтов, добавленных в $\mathcal{C}(t)$ за
        время наблюдения (в ходе разрешения конфликтов методом управления разъездом).
\end{itemize}

\subsection{Проверка статистической значимости}

Поскольку как порядок поступления заказов, так и начальное распределение агентов по графу разыгрываются 
псевдослучайно, каждая из метрик~(\ref{e:metric-orders})--(\ref{e:metric-speed}) при фиксированной 
конфигурации графа, числа агентов и метода разрешения конфликтов представляет собой случайную величину. 
Различие средних значений метрики между двумя методами, наблюдаемое на конечном числе запусков, может 
быть вызвано как действительным превосходством одного из них, так и случайным разбросом входных условий. 
Чтобы отделить систематический эффект от шума, а заодно убедиться в том, что наблюдаемые различия имеют 
практический смысл, а не только формально статистический, используем методы  проверки статистической 
значимости. Сравнение двух методов сводится при этом к сравнению двух независимых выборок, составленных 
из значений метрики, полученных при различных запусках симулятора.  Ниже описаны методы, которые будут 
применяться для такого сравнения.

В качестве основного критерия используется U-критерий Манна--Уитни~\cite{mann-test} 
(в реализации \texttt{scipy.stats.mannwhitneyu}). Нулевая гипотеза $H_0$ состоит в том, что 
распределения $X$ и $Y$ совпадают, альтернатива $H_1$ --- в том, что одна из выборок стохастически 
доминирует над другой. При $p < \alpha$ (используется уровень значимости $\alpha = 0.05$) гипотеза $H_0$ отвергается.

Выбор именно этого критерия обусловлен следующими соображениями. Во-первых, распределения исследуемых 
метрик заметно отличаются от нормального: $N(T)$ является дискретной счётной величиной, а $\bar\tau(T)$ и 
$\bar v(T)$ --- усреднениями по множеству маршрутов с тяжёлыми хвостами, возникающими при возникновении 
заторов. Это делает применение $t$-критерия Стьюдента некорректным. Во-вторых, критерий Манна--Уитни не 
требует предположений о форме распределения и устойчив к выбросам, что важно при сравнении методов, 
у которых отдельный неудачный прогон может дать аномально большое значение $\bar\tau(T)$. При множественных 
попарных сравнениях для контроля ложноположительных результатов используется пошаговая поправка Холма~\cite{holm} 
(в реализации \texttt{statsmodels.stats.multitest.multipletests}).

Чтобы дополнить $p$-value оценкой практической значимости, совместно с критерием Манна--Уитни 
вычисляется дельта Клиффа~\cite{cliff}. Для качественной интерпретации используется шкала 
Романо~\cite{romano}: эффект считается пренебрежимо малым при $|\delta| < 0{,}147$, малым при 
$|\delta| < 0{,}33$, средним при $|\delta| < 0{,}474$ и большим в остальных случаях.

Два метода считаются практически различимыми по выбранной метрике, если одновременно выполняются 
условия $p^{\mathrm{holm}} < 0{,}05$ и $|\delta| \geq 0{,}147$. Случай 
$p^{\mathrm{holm}} < 0{,}05$ при пренебрежимо малом $|\delta|$ интерпретируется как 
статистический артефакт большого объёма выборки и во внимание не принимается; случай 
$p^{\mathrm{holm}} \geq 0{,}05$ означает, что наблюдаемых данных недостаточно для отклонения 
гипотезы о совпадении распределений.

\subsection{Результаты работы}

Каждый запуск симулятора определяется четырьмя факторами: графом дорожной сети $G$, числом агентов 
$|A|$, методом разрешения конфликтов и используемым планировщиком маршрутов. Длительность симуляции 
ограничена 12 часами (43200 секунд) виртуального времени. Для набора экспериментальной выборки 
каждая конфигурация запускалась несколько раз (но не более двухсот) с суммарным ограничением по 
времени работы в 5 минут (10 минут для тяжелых конфигураций).

В сравнении участвуют:

\begin{itemize}
    \item графы трех размеров, указанные в раздел~\ref{sec:final_data};
    \item конфигурации агентов в размерах 10, 50, 100, 500, 1000 штук;
    \item методы разрешения конфликтов:
        \begin{itemize}
            \item управление разъездом (далее \enquote{разъезд}),
            \item глобальные ограничения через механизм семафоров
                (далее \enquote{семафор});
        \end{itemize}
    \item планировщики маршрутов:
        \begin{itemize}
            \item классический A*,
            \item адаптивный A*,
            \item CCBS (используется без механизма разрешения конфликтов).
        \end{itemize}
\end{itemize}

Поскольку абсолютные значения метрик~(\ref{e:metric-orders})--(\ref{e:metric-speed}) изменяются на 
порядки между парами $(G, |A|)$, объединение результатов симуляции из разных пар в одну выборку 
маскирует сравниваемый эффект под влияние масштаба. Поэтому каждая пара $(G, |A|)$ рассматривается 
как самостоятельная страта: все попарные сравнения и поправка Холма применяются только внутри 
страты, а итоговый вывод получается агрегированием результатов по всем доступным стратам.

Основным объектом сравнения является метод разрешения конфликтов, а планировщик играет вспомогательную 
роль. В связи с этим анализ результатов разбит на три шага:

\begin{enumerate}
    \item для каждого метода разрешения конфликтов внутри страты выбирается
        планировщик, с которым этот метод показывает лучший результат;
    \item методы разрешения конфликтов сравниваются между собой в их
        наилучших конфигурациях;
    \item полученные конфигурации сравниваются с CCBS.
\end{enumerate}

\subsubsection{Выбор метода планирования для методов разрешения конфликтов}

Для каждого метода разрешения конфликтов $R$ и каждой страты $(G, |A|)$ сравниваются две выборки:
$R$ в паре с классическим A* против $R$ в паре с адаптивным A*, по каждой из трех основных 
метрик~(\ref{e:metric-orders})--(\ref{e:metric-speed}).

По каждой метрике голос отдается в пользу того планировщика, для которого метрика лучше в нужном 
направлении (число заказов и средняя скорость --- больше, время заказа --- меньше), при условии что 
$p^{\mathrm{holm}} < 0{,}05$ и $|\delta| \geq 0{,}147$. В остальных случаях голос засчитывается как 
ничейный. Побеждает планировщик, получивший хотя бы два голоса из трех; при делении $1{-}1{-}1$ 
приоритет имеет метрика $N(T)$; при трех ничьих победитель не определен.

В таблице~\ref{tab:router-reverse} представлен фрагмент сравнения методов планирования для метода \enquote{разъезд}
по метрике средней скорости.

\begin{table}
\caption{Сравнение методов планирования для метода \enquote{разъезд} по метрике средней скорости}
\label{tab:router-reverse}
\begin{tabular}{|l|c|c|c|c|c|c|}
\hline
\multirow{2}{*}{Граф} & \multirow{2}{*}{\makecell{Число\\агентов}} & \multicolumn{2}{c|}{медиана} & \multirow{2}{*}{$\delta$} & \multirow{2}{*}{$p^{\mathrm{holm}}$} & \multirow{2}{*}{результат} \\
\cline{3-4}
 & & A* & адапт. A* & & & \\
\hline
\multirow{3}{*}{мал.} & $  10$ & $2.063$ & $2.069$ & $-0.927$ & $9.882\cdot{-39}$ & адапт. A* \\
                      & $ 100$ & $1.934$ & $1.993$ & $-1     $ & $1.574\cdot{-28}$ & адапт. A* \\
                      & $1000$ & $1.575$ & $1.944$ & $-1     $ & $7.589\cdot{-8}$ & адапт. A* \\
\hline
\multirow{3}{*}{ср.}  & $  10$ & $2.137$ & $2.142$ & $-0.893$ & $5.284\cdot{-36}$ & адапт. A* \\
                      & $ 100$ & $2.099$ & $2.117$ & $-1     $ & $2.583\cdot{-24}$ & адапт. A* \\
                      & $1000$ & $1.976$ & $2.072$ & $-1     $ & $0.00199   $ & адапт. A* \\
\hline
\multirow{3}{*}{бол.} & $  10$ & $2.116$ & $2.138$ & $-0.999$ & $9.198\cdot{-45}$ & адапт. A* \\
                      & $ 100$ & $2.093 $ & $2.146$ & $-1     $ & $4.232\cdot{-24} $ & адапт. A* \\
                      & $1000$ & $2.031$ & $2.146$ & $-1     $ & $0.0035  $ & адапт. A* \\
\hline
\end{tabular}
\end{table}

Из таблицы видно, что при небольшом числе агентов вклад метода планирования практически незаметен, что обусловлено
меньшим числом конфликтов в таких конфигурациях и соответственно меньшим эффектом от разрешения конфликтов в целом.
При увеличении числа агентов начинают играть б\'ольшую роль физические ограничения локации, а также меньший объем данных,
который был собран в таких запусках. Адаптивный A* показывает себя лучше, чем неадаптивный.

В таблице~\ref{tab:router-sem} представлено сравнение методов планирования для метода \enquote{семафор}
по метрике числа заказов. Выводы из таблицы можно сделать аналогичные.

\begin{table}
\caption{Сравнение методов планирования для метода \enquote{семафор} по метрике числа заказов}
\label{tab:router-sem}
\begin{tabular}{|l|c|c|c|c|c|c|}
\hline
\multirow{2}{*}{Граф} & \multirow{2}{*}{\makecell{Число\\агентов}} & \multicolumn{2}{c|}{медиана} & \multirow{2}{*}{$\delta$} & \multirow{2}{*}{$p^{\mathrm{holm}}$} & \multirow{2}{*}{результат} \\
\cline{3-4}
 & & A* & адапт. A* & & & \\
\hline
\multirow{3}{*}{мал.} & $  10$ & $ 2064.5$ & $ 2067  $ & $-0.0574  $ & $0.836   $ & ничья  \\
                      & $ 100$ & $18081  $ & $18656  $ & $-1       $ & $6.507\cdot10^{-30}$ & адапт. A* \\
                      & $1000$ & $66273  $ & $89888  $ & $-1       $ & $1.991\cdot10^{-11}$ & адапт. A* \\
\hline
\multirow{3}{*}{ср.}  & $  10$ & $  606  $ & $  608.5$ & $-0.07365 $ & $0.585   $ & ничья  \\
                      & $ 100$ & $ 5942  $ & $ 5973.5$ & $-0.508122$ & $7.366\cdot10^{-6}$ & адапт. A* \\
                      & $1000$ & $51529  $ & $54821  $ & $-1       $ & $6.625\cdot10^{-5}$ & адапт. A* \\
\hline
\multirow{3}{*}{бол.} & $  10$ & $  309.5$ & $  312  $ & $-0.1913  $ & $0.0138  $ & адапт. A* \\
                      & $ 100$ & $ 3115  $ & $ 3184  $ & $-0.975684$ & $4.289\cdot10^{-18}$ & адапт. A* \\
                      & $1000$ & $28104  $ & $31569  $ & $-1       $ & $0.00024$ & адапт. A* \\
\hline
\end{tabular}
\end{table}

При агрегировании результата по всем метрикам получаются следующие результаты, которые представлены
на рисунке~\ref{fig:router-final} (цвет ячеек отражает дельту Клиффа по метрике средней скорости, 
статистическая значимость выражена звездочками).

\begin{figure}
\includegraphics[width=\textwidth]{astar_comparison}
\caption{Итоговые результаты выбора метода планирования для методов разрешения конфликтов}
\label{fig:router-final}
\end{figure}

Таким образом, при небольшом числе агентов может быть использован любой метода планирования, но в целом
адаптивный A* показывает себя лучше в большинстве сценариев и ни в одном сценарии не показывает себя хуже.
Для всех методов разрешения конфликтов выбираем адаптивный A*.

\subsubsection{Сравнение методов разрешения конфликтов}

В каждой страте $(G, |A|)$ метод \enquote{разъезд} в выбранной для него конфигурации сравнивается с 
методом \enquote{семафор} в выбранной для него конфигурации.

Каждой паре \enquote{метод~$\times$ метрика} внутри страты сопоставляется голос по тому же правилу, 
что и в предыдущем шаге. Победитель страты определяется большинством голосов по всем 
метрикам и всем парам разложения.

В таблице~\ref{tab:resolvers} представлено сравнение методов разрешения конфликтов по метрике среднего времени заказа.

\begin{table}
\caption{Сравнение методов разрешения конфликтов по метрике среднего времени заказа}
\label{tab:resolvers}
\begin{tabular}{|l|c|c|c|c|c|c|}
\hline
\multirow{2}{*}{Граф} & \multirow{2}{*}{\makecell{Число\\агентов}} & \multicolumn{2}{c|}{медиана} & \multirow{2}{*}{$\delta$} & \multirow{2}{*}{$p^{\mathrm{holm}}$} & \multirow{2}{*}{результат} \\
\cline{3-4}
 & & \enquote{семафор} & \enquote{разъезд} & & & \\
\hline
\multirow{5}{*}{мал.} & $  10$ & $107  $ & $104.4$ &  $0.964 $ & $7.112\cdot10^{-32}$ & \enquote{семафор} \\
                      & $  50$ & $113.5$ & $110.3$ &  $1     $ & $7.64\cdot10^{-34} $ & \enquote{семафор} \\
                      & $ 100$ & $117  $ & $115.8$ &  $0.946$ & $5.505\cdot10^{-25}$ & \enquote{семафор} \\
                      & $ 500$ & $146.9$ & $176.9$ & -$1     $ & $8.479\cdot10^{-14}$ & \enquote{разъезд}   \\
                      & $1000$ & $168  $ & $239.3$ & -$1     $ & $8.392\cdot10^{-07}$ & \enquote{разъезд}   \\
 \hline
 \multirow{5}{*}{ср.} & $  10$ & $355.2$ & $354.7$ &  $0.063$ & $0.8903   $ & ничья       \\
                      & $  50$ & $363.3$ & $359.5$ &  $0.619$ & $7.847\cdot10^{-14}$ & \enquote{семафор} \\
                      & $ 100$ & $366.4$ & $361.2$ &  $0.932$ & $1.244\cdot10^{-18}$ & \enquote{семафор} \\
                      & $ 500$ & $381.7$ & $376.9$ &  $1     $ & $1.22\cdot10^{-08} $ & \enquote{семафор} \\
                      & $1000$ & $393.3$ & $393.6$ & -$0.367$ & $0.527    $ & ничья       \\
\hline
\multirow{5}{*}{бол.} & $  10$ & $687.3$ & $690.7$ & -$0.108$ & $0.5672   $ & ничья       \\
                      & $  50$ & $677.4$ & $678.2$ & -$0.05$ & $1        $ & ничья       \\
                      & $ 100$ & $680.3$ & $677  $ &  $0.31$ & $0.00363 $ & \enquote{семафор} \\
                      & $ 500$ & $684.6$ & $681.3$ &  $0.721$ & $2.382\cdot10^{-05}$ & \enquote{семафор} \\
                      & $1000$ & $687.5$ & $682.8$ &  $0.873 $ & $0.00419 $ & \enquote{семафор} \\
\hline
\end{tabular}
\end{table}

Итоговые результаты по всем метрикам представлены на рисунке~\ref{fig:resolvers-final} 
(цвет ячеек отражает дельту Клиффа по метрике средней скорости).

\begin{figure}
\includegraphics[width=\textwidth]{resolver_comparison}
\caption{Итоговые результаты выбора метода разрешения конфликтов}
\label{fig:resolvers-final}
\end{figure}

По результатам этого сравнения видно, что при небольшом числе агентов относительно размера локации
выбор метода разрешения конфликтов не имеет большой разницы из-за малого числа самих конфликтов.
В большинстве сценариев лучше себя показывает метод \enquote{семафор}, за исключением случаев с большим числом
агентов на меньших графах, в которых он уступает методу разъездов.

\subsubsection{Сравнение с CCBS}
\label{sec:res-ccbs}

Для каждой страты $(G, |A|)$ с $|A| \in \{10, 100, 1000\}$ строятся две пары сравнения: CCBS против 
\enquote{разъезда} в его наилучшей конфигурации и CCBS против \enquote{семафора} в его наилучшей 
конфигурации, все шесть тестов одной страты объединяются в одно семейство для поправки Холма.

В таблице~\ref{tab:ccbs} представлен небольшой фрагмент сравнения методов разрешения конфликтов с CCBS.

\begin{table}
\caption{Сравнение методов разрешения конфликтов с CCBS (фрагмент)}
\label{tab:ccbs}
\begin{tabular}{|c|c|c|c|c|c|c|c|c|}
\hline
Граф & $|A|$ & метод & метрика & $\delta$ & $p^{\mathrm{holm}}$ & результат \\
\hline
\multirow{2}{*}{мал.} &  \multirow{2}{*}{10}  & \enquote{разъезд}    & число заказов    & $-0.898 $ & $1.466\cdot10^{-36}$ & разрешатель \\
&                       & \enquote{семафор}    & среднее время    & $ 1      $ & $1.464\cdot10^{-44}$ & разрешатель \\
\hline
\multirow{3}{*}{ср.}  &  \multirow{3}{*}{100} & \enquote{разъезд}    & число заказов    & $-0.089$ & $1        $ & ничья      \\
&                       & \enquote{разъезд}    & среднее время    & $ 0.079$ & $1        $ & ничья      \\
&                       & \enquote{семафор}    & средняя скорость & $-1      $ & $5.318\cdot10^{-11}$ &  разрешатель \\
\hline
\multirow{3}{*}{бол.} & \multirow{3}{*}{1000} & \enquote{разъезд}    & число заказов    & $ 0.534 $ & $0.1385   $ & ничья      \\
&                       & \enquote{семафор}    & среднее время    & $ 0.576 $ & $0.1336   $ & ничья      \\
&                       & \enquote{семафор}    & средняя скорость & $ 0.919 $ & $0.003145 $ & CCBS      \\
\hline
\end{tabular}
\end{table}

Итоговые результаты по всем метрикам представлены на рисунке~\ref{fig:ccbs-final} 
(цвет ячеек отражает дельту Клиффа по метрике числа заказов).

\begin{figure}
\includegraphics[width=\textwidth]{ccbs_comparison}
\caption{Итоговые результаты сравнения методов разрешения конфликтов с CCBS}
\label{fig:ccbs-final}
\end{figure}

По результатам итогового сравнения видно, что методы разрешения конфликтов показывают себя не хуже, чем
метод многоагентного планирования. В большинстве случаев эти методы превосходят CCBS, но начинают уступать
при увеличении размера задачи. Однако стоит отметить, что при увеличении размера задачи начинает испытывать трудности и метод CCBS,
что выражается в его скорости работы и сходимости.

Основные выводы из анализа результатов такие, что методы локального разрешения конфликтов могут быть использованы
для управления автономными агентами и не уступают методам многоагентного планирования по качеству.
Вместе с тем, методам разрешения конфликтов очень помогает использование информации о состоянии среды, 
что может быть выражено в использовании адаптивного одногоагентного метода построения маршрутов.
